%%%%%%%%%%%%
%
% $Autor: Wings $
% $Datum: 2019-03-05 08:03:15Z $
% $Pfad: ListOfParts.tex $
% $Version: 4250 $
% !TeX spellcheck = en_GB/de_DE
% !TeX encoding = utf8
% !TeX root = manual 
% !TeX TXS-program:bibliography = txs:///biber
%
%%%%%%%%%%%%

\chapter{Liste der Teile}

In diesem Kapitel werden die Bestandteile des Tischförderbands aufgeführt, welche für die Inbetriebnahme benötigt werden. Falls Komponenten beschädigt sein sollten oder fehlen, kontaktieren Sie uns bitte umgehend.

\bigskip
\bigskip

% Feste Höhe der Bildzelle
\newlength{\BildZellenHoehe}
\setlength{\BildZellenHoehe}{3cm}

\begin{table}
	\caption{Liste der im Lieferumfang enthaltenen Teile}
	\begin{center}
		\renewcommand{\arraystretch}{1.2}
		
		\begin{tabular}{|l|l|>{\centering\arraybackslash}m{6cm}|c|}
			\hline
			Nr. & Bezeichnung & Bild & Anzahl \\
			\hline  
			
			1 
			& Tischförderband
			& \parbox[c][\BildZellenHoehe][c]{\linewidth}{\centering
				\includegraphics[
				height=0.9\BildZellenHoehe,
				keepaspectratio
				]{NewFiles/images/Bild_Förderband_Handbuch.png}}
			& 1 \\
			\hline
			
			2
			& Netzkabel mit Netzteil
			& \parbox[c][\BildZellenHoehe][c]{\linewidth}{\centering
				\includegraphics[
				height=0.9\BildZellenHoehe,
				keepaspectratio
				]{NewFiles/images/Bild_Netzkabel_mit_Netzteil_Handbuch.jpg}}
			& 1 \\
			\hline
			
			3
			& Benutzerhandbuch
			& \parbox[c][\BildZellenHoehe][c]{\linewidth}{\centering
				
\includegraphics[
				height=0.9\BildZellenHoehe,
				keepaspectratio
				]{NewFiles/images/Bild_Benutzerhandbuch_Handbuch.png}}
			& 1 \\
			\hline
			
			4
			& Transportboxen
			& \parbox[c][\BildZellenHoehe][c]{\linewidth}{\centering
				
\includegraphics[
				height=0.9\BildZellenHoehe,
				keepaspectratio
				]{NewFiles/images/Bild_Transportboxen_Handbuch.jpg}}
			& 3 \\
			\hline
			
			5
			& Hilfswerkzeug
			& \parbox[c][\BildZellenHoehe][c]{\linewidth}{\centering
				\includegraphics[
				height=0.9\BildZellenHoehe,
				keepaspectratio
				]{NewFiles/images/Bild_Hilfswerkzeug_Handbuch.jpg}}
			& 1 \\
			\hline
			
		\end{tabular}
	\end{center}
\end{table}